% GENERATED FILE. Edit canonical lessons, then run npm run generate:books.
% canonical-source-hash: fnv1a64:28917671a4c433d4
% canonical-lessons: GU-C285-cor, GU-C285-cori, GU-C285-nasvu, GU-C285-pakadvu, GU-C285-loko

\chapter{A Thief in the Night}
\label{ch:gu-a-thief-in-the-night}
\hlchaptermodality{285}

\begin{chapteropening}
\textbf{By the end of this chapter:} \emph{I can say that a thief came in the night and what happened next.}

Complete the last lesson of chapter 285: I can say that a thief came in the night and what happened next.
\end{chapteropening}

\section[\texorpdfstring{\gu{ચોર} (cor)}{ચોર (cor)}]{\gu{ચોર} (cor) — a thief}

\label{lesson:GU-C285-cor}



\begin{quote}
Before the new one: say the Gujarati for to pester, then the Gujarati for to hope.
\end{quote}

\subsection*{You'll want to know: \gu{ચોર}}
\textbf{\gu{ચોર}} — \emph{cor} — \textquotedblleft{}a thief\textquotedblright{}.

\begin{quote}
\gu{ગઈકાલે} \gu{રાત્રે} \gu{એક} \gu{ચોર} \gu{અમારા} \gu{ઘરમાં} \gu{આવ્યો}.
\end{quote}

\begin{quote}
\emph{gaīkāle rātre ek cor amārā gharmā̃ āvyo.}
\end{quote}

\begin{quote}
Last night a thief got into our house.
\end{quote}

A masculine or feminine noun, by the person meant.

\begin{cousinweb}
\textbf{\gu{ચોર}} (\emph{cor}) is inherited from Sanskrit \emph{cōrá-}, \textquotedblleft{}a thief\textquotedblright{}, and CDIAL lists the Gujarati word (CDIAL 4931). The verb \textbf{\gu{ચોરવું}}, to steal, belongs to the same family.
\end{cousinweb}

\begin{grammarlens}[title={What you've built}]
One more word for telling what happened.
\end{grammarlens}

\subsection*{Your turn}
\begin{itemize}
\raggedright
  \item \emph{Say it:} \emph{cor}
  \item \emph{Say it:} \emph{cor}, once more
  \item \emph{Say it:} \emph{gaīkāle rātre ek cor amārā gharmā̃ āvyo.} — Last night a thief got into our house.
  \item \emph{Recall:} say the Gujarati for to pester, then the Gujarati for to hope, then say \emph{cor} again
\end{itemize}

\begin{tcolorbox}[breakable,colback=teal!4,colframe=teal!35!black,title={Before you move on}]
What does \gu{ચોર} mean? (\textquotedblleft{}A thief\textquotedblright{}.) Say it once more.
\end{tcolorbox}
\section[\texorpdfstring{\gu{ચોરી} (corī)}{ચોરી (corī)}]{\gu{ચોરી} (corī) — a theft}

\label{lesson:GU-C285-cori}



\begin{quote}
Before the new one: say the Gujarati for to hope, then the Gujarati for a thief.
\end{quote}

\subsection*{You'll want to know: \gu{ચોરી}}
\textbf{\gu{ચોરી}} — \emph{corī} — \textquotedblleft{}a theft\textquotedblright{}.

\begin{quote}
\gu{ગઈકાલે} \gu{રાત્રે} \gu{અમારી} \gu{ગલીમાં} \gu{ચોરી} \gu{થઈ}.
\end{quote}

\begin{quote}
\emph{gaīkāle rātre amārī galīmā̃ corī thaī.}
\end{quote}

\begin{quote}
There was a theft in our lane last night.
\end{quote}

A feminine noun.

\begin{cousinweb}
is inherited from Sanskrit \emph{caurikā}, \textquotedblleft{}theft\textquotedblright{} (CDIAL 4937). \textbf{\gu{ચોર}} is the person and \textbf{\gu{ચોરી}} what the person does; a theft takes place with \textbf{\gu{થવું}}: \textbf{\gu{ચોરી} \gu{થઈ}}.
\end{cousinweb}

\begin{grammarlens}[title={What you've built}]
One more word for telling what happened.
\end{grammarlens}

\subsection*{Your turn}
\begin{itemize}
\raggedright
  \item \emph{Say it:} \emph{corī}
  \item \emph{Say it:} \emph{corī}, once more
  \item \emph{Say it:} \emph{gaīkāle rātre amārī galīmā̃ corī thaī.} — There was a theft in our lane last night.
  \item \emph{Recall:} say the Gujarati for to hope, then the Gujarati for a thief, then say \emph{corī} again
\end{itemize}

\begin{tcolorbox}[breakable,colback=teal!4,colframe=teal!35!black,title={Before you move on}]
What does \gu{ચોરી} mean? (\textquotedblleft{}A theft\textquotedblright{}.) Say it once more.
\end{tcolorbox}
\section[\texorpdfstring{\gu{નાસવું} (nāsvũ)}{નાસવું (nāsvũ)}]{\gu{નાસવું} (nāsvũ) — to flee}

\label{lesson:GU-C285-nasvu}



\begin{quote}
Before the new one: say the Gujarati for a thief, then the Gujarati for a theft.
\end{quote}

\subsection*{You'll want to know: \gu{નાસવું}}
\textbf{\gu{નાસવું}} — \emph{nāsvũ} — \textquotedblleft{}to flee\textquotedblright{}.

\begin{quote}
\gu{ચોરે} \gu{પાકીટ} \gu{ચોર્યું} \gu{અને} \gu{તે} \gu{નાસી} \gu{ગયો}.
\end{quote}

\begin{quote}
\emph{core pākīṭ coryũ ane te nāsī gayo.}
\end{quote}

\begin{quote}
The thief stole the wallet and ran off.
\end{quote}

Past (he …): \gu{તે} \gu{ચોર} \gu{નાસી} \gu{ગયો}, \emph{te cor nāsī gayo}.

\begin{cousinweb}
\textbf{\gu{નાસવું}} (\emph{nāsvũ}) is inherited from Sanskrit \emph{náśyati}, \textquotedblleft{}is lost, runs away\textquotedblright{}, and CDIAL glosses the Gujarati verb \textquotedblleft{}to flee\textquotedblright{} (CDIAL 7027). It is usually said with \textbf{\gu{જવું}}: \textbf{\gu{નાસી} \gu{ગયો}}, he ran off.
\end{cousinweb}

\begin{grammarlens}[title={What you've built}]
One more word for telling what happened.
\end{grammarlens}

\subsection*{Your turn}
\begin{itemize}
\raggedright
  \item \emph{Say it:} \emph{nāsvũ}
  \item \emph{Say it:} \emph{nāsvũ}, once more
  \item \emph{Say it:} \emph{core pākīṭ coryũ ane te nāsī gayo.} — The thief stole the wallet and ran off.
  \item \emph{Recall:} say the Gujarati for a thief, then the Gujarati for a theft, then say \emph{nāsvũ} again
\end{itemize}

\begin{tcolorbox}[breakable,colback=teal!4,colframe=teal!35!black,title={Before you move on}]
What does \gu{નાસવું} mean? (\textquotedblleft{}To flee\textquotedblright{}.) Say it once more.
\end{tcolorbox}
\section[\texorpdfstring{\gu{પકડવું} (pakaḍvũ)}{પકડવું (pakaḍvũ)}]{\gu{પકડવું} (pakaḍvũ) — to catch}

\label{lesson:GU-C285-pakadvu}



\begin{quote}
Before the new one: say the Gujarati for a theft, then the Gujarati for to flee.
\end{quote}

\subsection*{You'll want to know: \gu{પકડવું}}
\textbf{\gu{પકડવું}} — \emph{pakaḍvũ} — \textquotedblleft{}to catch\textquotedblright{}.

\begin{quote}
\gu{પોલીસે} \gu{ચોરને} \gu{બજારમાં} \gu{પકડ્યો}.
\end{quote}

\begin{quote}
\emph{polīse corne bajārmā̃ pakaḍyo.}
\end{quote}

\begin{quote}
The police caught the thief in the market.
\end{quote}

Past, with \gu{તેણે} (he or she …): \gu{તેણે} \gu{ચોરને} \gu{પકડ્યો}, \emph{teṇe corne pakaḍyo}.

\begin{cousinweb}
\textbf{\gu{પકડવું}} (\emph{pakaḍvũ}) comes from \emph{*pakkaḍ}, \textquotedblleft{}to seize\textquotedblright{}, a form CDIAL reconstructs from its descendants, and CDIAL lists the Gujarati verb (CDIAL 7619). A person caught takes \textbf{\gu{ને}}, and the verb agrees with him: \textbf{\gu{ચોરને} \gu{પકડ્યો}}.
\end{cousinweb}

\begin{grammarlens}[title={What you've built}]
One more word for telling what happened.
\end{grammarlens}

\subsection*{Your turn}
\begin{itemize}
\raggedright
  \item \emph{Say it:} \emph{pakaḍvũ}
  \item \emph{Say it:} \emph{pakaḍvũ}, once more
  \item \emph{Say it:} \emph{polīse corne bajārmā̃ pakaḍyo.} — The police caught the thief in the market.
  \item \emph{Recall:} say the Gujarati for a theft, then the Gujarati for to flee, then say \emph{pakaḍvũ} again
\end{itemize}

\begin{tcolorbox}[breakable,colback=teal!4,colframe=teal!35!black,title={Before you move on}]
What does \gu{પકડવું} mean? (\textquotedblleft{}To catch\textquotedblright{}.) Say it once more.
\end{tcolorbox}
\section[\texorpdfstring{\gu{લોકો} (loko)}{લોકો (loko)}]{\gu{લોકો} (loko) — people}

\label{lesson:GU-C285-loko}



\begin{quote}
Before the new one: say the Gujarati for to flee, then the Gujarati for to catch.
\end{quote}

\subsection*{You'll want to know: \gu{લોકો}}
\textbf{\gu{લોકો}} — \emph{loko} — \textquotedblleft{}people\textquotedblright{}.

\begin{quote}
\gu{પોલીસનું} \gu{કામ} \gu{લોકો} \gu{પર} \gu{પણ} \gu{આધાર} \gu{રાખે} \gu{છે}.
\end{quote}

\begin{quote}
\emph{polīsnũ kām loko par paṇ ādhār rākhe che.}
\end{quote}

\begin{quote}
The work of the police depends on people too.
\end{quote}

A masculine noun, used in the plural.

\begin{cousinweb}
\textbf{\gu{લોકો}} (\emph{loko}) is the plural of \emph{lok}, from Sanskrit \emph{lōká-}, \textquotedblleft{}open space, the world\textquotedblright{}, which later also meant \textquotedblleft{}people\textquotedblright{} (CDIAL 11119). Gujarati uses it in the plural, and a verb with it is plural too: \textbf{\gu{લોકો} \gu{આવ્યા}}.
\end{cousinweb}

\begin{grammarlens}[title={What you've built}]
One more word for telling what happened.
\end{grammarlens}

\subsection*{Your turn}
\begin{itemize}
\raggedright
  \item \emph{Say it:} \emph{loko}
  \item \emph{Say it:} \emph{loko}, once more
  \item \emph{Say it:} \emph{polīsnũ kām loko par paṇ ādhār rākhe che.} — The work of the police depends on people too.
  \item \emph{Recall:} say the Gujarati for to flee, then the Gujarati for to catch, then say \emph{loko} again
\end{itemize}

\begin{tcolorbox}[breakable,colback=teal!4,colframe=teal!35!black,title={Before you move on}]
What does \gu{લોકો} mean? (\textquotedblleft{}People\textquotedblright{}.) Say it once more.
\end{tcolorbox}
